It is region $R_1$ that will dominate the value of the integral $I(\phi)$.  We
will first perform a Taylor series expansion of all the terms around their
sample versions, and then take probability limits of the resulting expressions.

First, we will show that integrals over $\normhat$ in $R_1$ are sufficiently
close to the corresponding Gaussian integrals over $\rdom{\thetadim}$ for
sufficiently large $\delta_1$.

%%%%%%%%%%%%%%%%%%%%%%%%%%%%%%%%%%%%%%%%%%%%%%%%%%%%%%%%%%%%%%%%%%%%%%%%%
%%%%%%%%%%%%%%%%%%%%%%%%%%%%%%%%%%%%%%%%%%%%%%%%%%%%%%%%%%%%%%%%%%%%%%%%%
%%%%%%%%%%%%%%%%%%%%%%%%%%%%%%%%%%%%%%%%%%%%%%%%%%%%%%%%%%%%%%%%%%%%%%%%%

\begin{lem}\lemlabel{region_r1_gaussian_moments}
%
Let \assuref{core_bclt_assu} hold.  Let $\tau^k$ represent the
$\thetadim$-dimensional array of outer products of $\tau$.  When $\regevent$ occurs and
$\delta_1 \ge \sqrt{24 / \infoev}$ then, for any $k \ge 0$,
%
\begin{align*}
%
\int_{R_1} \normhatarg{\tau} \tau^k d\tau =
    \expect{\normhat}{\tau^k} + O(N^{-3}).
%
\end{align*}
%
\begin{proof}
%
By definition,
%
\begin{align*}
%
\int_{R_1} \normhatarg{\tau} \tau^k d\tau  - \expect{\normhat}{\tau^k}
={}&
\int_{R_2 \bigcup R_3} \normhatarg{\tau} \tau^k d\tau
%
\end{align*}
%
We will now use a technique similar to that of the proof of
\lemref{region_r2} to control the norm of the final integral.
Recall that, on $\regevent$, the minimum eigenvalue of $\infohat$, $\infoevhat$,
is at least $\frac{1}{2} \infoev$ (\lemref{regular_event} \eventref{infoev}),
so that $\normalizer^{-1} = O(1)$.
%
\begin{align*}
%
\MoveEqLeft
\norm{
    \int_{R_2 \bigcup R_3}
            \normhatarg{\tau} \tau^k d\tau
}_2
\\\le&
    \int_{R_2 \bigcup R_3}
            \normhatarg{\tau} \norm{\tau}^k_2 d\tau
\\={}&
\normalizer^{-1} \int_{R_2 \bigcup R_3}
        \exp\left(-\frac{1}{2} \infohat \tau^2 \right)
            \norm{\tau}^k_2 d\tau
\\\le&
\normalizer^{-1} \int_{R_2 \bigcup R_3}
        \exp\left(-\frac{1}{4} \infoev \norm{\tau}_2^2 \right)
            \norm{\tau}^k_2 d\tau
\\\le&
\sup_{\tau \in R_2 \bigcup R_3}
    \exp\left(-\frac{1}{8} \infoev \norm{\tau}_2^2 \right)
\normalizer^{-1} \int_{R_2 \bigcup R_3}
        \exp\left(-\frac{1}{8} \infoev \norm{\tau}_2^2 \right)
            \norm{\tau}^k_2 d\tau
\\={}&
\exp\left(-\frac{1}{8} \infoev \delta_1^2 (\log N)^2 \right)
\normalizer^{-1} \int_{R_2 \bigcup R_3}
        \exp\left(-\frac{1}{8} \infoev \norm{\tau}_2^2 \right)
            \norm{\tau}^k_2 d\tau
\\={}& O(N^{-3}),
%
\end{align*}
%
where the final line takes $\delta_1^2 \ge \frac{24}{\infoev}$, that $\log N <
(\log N)^2$ for $N \ge 3$, that $\normalizer^{-1} = O(1)$, and that normal
moments are finite.
%
\end{proof}
%
\end{lem}

%%%%%%%%%%%%%%%%%%%%%%%%%%%%%%%%%%%%%%%%%%
%%%%%%%%%%%%%%%%%%%%%%%%%%%%%%%%%%%%%%%%%%
%%%%%%%%%%%%%%%%%%%%%%%%%%%%%%%%%%%%%%%%%%

Next, we series expand the likelihood term.

% Note that you need a standalone lemma for the log likelihood expansion
% to deal with the normalizing constant.
\begin{lem}\lemlabel{region_r1_loglik_expansion}
%
Let \assuref{core_bclt_assu} hold.  For any $\xvec$ such that $\regevent$
occurs, and for $\theta \in R_1$,
%
\begin{align*}
%
\MoveEqLeft
\exp\left(N \likhat(\theta) - N \likhat(\thetahat)\right) = \\
& \exp\left(\frac{1}{2} \infohat \tau^2 \right) \times
\left( 1 + N^{-1/2} \frac{1}{6} \likhatk{3}(\theta) \tau^3 +
    \ordlog{N^{-1}} \right).
%
\end{align*}
%
\begin{proof}
%

Using \eqref{taylor_expansion}, write the log likelihood term as
%
\begin{align}
%
\MoveEqLeft
\exp\left(N \likhat(\theta) - N \likhat(\thetahat)\right) =  \nonumber \\
& \exp\left(\frac{1}{2} \likhatk{2}(\thetahat) \tau^2 \right) \times \nonumber \\
& \exp\left(
        N^{-1/2} \frac{1}{6} \likhatk{3}(\thetahat) \tau^3 +
        N^{-1} \frac{1}{6} \tresid{4}{\likhat(\theta), \thetahat, \theta} \tau^4
    \right). \eqlabel{r1_loglike_term_factorization}
%
\end{align}
%
We will expand the second term of \eqref{r1_loglike_term_factorization} using
the expansion of the  exponential function around 0, which holds for all
$z$:
%
\begin{align*}
\exp(z) ={}& 1 + z + z^2 + \frac{1}{2} \int_0^1 (1 - t)^2 \exp(tz) dt z^3 \\
={}& 1 + z + O(|z|^2).
\end{align*}
%
To match \eqref{r1_loglike_term_factorization}, we need to evaluate at
%
\begin{align*}
%
z = N^{-1/2} \frac{1}{6} \likhatk{3}(\thetahat) \tau^3 +
    N^{-1} \tresid{4}{\likhat(\theta), \thetahat, \theta} \tau^4.
%
\end{align*}
%
When $\regevent$ occurs (\lemref{regular_event} \eventref{ulln}), we
have
%
\begin{align*}
%
\norm{N^{-1} \tresid{4}{\likhat(\theta), \thetahat, \theta} \tau^4}_2 \le&
N^{-1} \sup_{\theta \in R_1}
    \norm{\tresid{4}{\likhat(\theta), \thetahat, \theta}}_2
    \sup_{\tau \in R_1}\norm{\tau^3}_2 \\
={}& \ordlog{N^{-1}}\\
%
\norm{N^{-1/2} \frac{1}{6} \likhatk{3}(\thetahat) \tau^3}_2 \le&
N^{-1/2} \frac{1}{6} \norm{\likhatk{3}(\thetahat)}_2
    \sup_{\tau \in R_1}\norm{\tau^3}_2\\
={}& \ordlog{N^{-1/2}} \Rightarrow\\
|z| ={}& \ordlog{N^{-1/2}},
%
\end{align*}
%
so that
%
\begin{align}
%
\exp\left(
        N^{-1/2} \frac{1}{6} \likhatk{3}(\thetahat) \tau^3 +
        N^{-1} \frac{1}{6} \tresid{4}{\likhat(\theta), \thetahat, \theta} \tau^4
    \right) = \nonumber\\
1 + N^{-1/2} \frac{1}{6} \likhatk{3}(\theta) \tau^3 +
    \ordlog{N^{-1}}
. \eqlabel{r1_exp_expansion}
%
\end{align}
%
Combining \eqref{r1_loglike_term_factorization, r1_exp_expansion} and
recognizing the definition of $\infohat$ gives the desired result.
%
\end{proof}
%
\end{lem}


%%%%%%%%%%%%%%%%%%%%%%%%%%%%%%%%%%%%%%%%%%%%%%%%%%%%%%%%%%%%
%%%%%%%%%%%%%%%%%%%%%%%%%%%%%%%%%%%%%%%%%%%%%%%%%%%%%%%%%%%%
%%%%%%%%%%%%%%%%%%%%%%%%%%%%%%%%%%%%%%%%%%%%%%%%%%%%%%%%%%%%

\begin{lem}\lemlabel{r1_normalizer_expansion}
%
Let \assuref{core_bclt_assu} hold.  When $\regevent$ occurs and $\delta_1$ is as
given in \lemref{region_r1_gaussian_moments},
%
\begin{align*}
%
\int_{R_1} \exp\left(
    N \likhat(\theta) - N \likhat(\thetahat)
\right) d\tau = \normalizer + \ordlog{N^{-1}}.
%
\end{align*}
%
\begin{proof}
%
Recall that, by definition,
%
\begin{align*}
%
\normalizer = \int \exp\left(-\frac{1}{2} \infohat \tau^2 \right) d\tau.
%
\end{align*}
%
Expanding using \lemref{region_r1_loglik_expansion} gives
%
\begin{align*}
%
\MoveEqLeft
\left| \int_{R_1} \exp\left(
    N \likhat(\theta) - N \likhat(\thetahat)
\right) d\tau - \normalizer \right|
\\={}&
\left|
\int_{R_1} \exp\left(\frac{1}{2} \infohat \tau^2 \right)
\left( 1 + N^{-1/2} \frac{1}{6} \likhatk{3}(\theta) \tau^3 +
    \ordlog{N^{-1}} \right) d\tau -
\int \exp\left(-\frac{1}{2} \infohat \tau^2 \right) d\tau
\right|
%
\\={}&
\left|
\int_{R_1} \exp\left(\frac{1}{2} \infohat \tau^2 \right)
\left( N^{-1/2} \frac{1}{6} \likhatk{3}(\theta) \tau^3 +
    \ordlog{N^{-1}} \right) d\tau -
\int_{R_2 \bigcup R_3} \exp\left(-\frac{1}{2} \infohat \tau^2 \right) d\tau
\right|
%
\\\le&
\frac{1}{6} N^{-1/2} \sup_{\theta \in R_1} \norm{
    \likhatk{3}(\theta) }_2
\expect{\normhat}{\ind{\tau \in R_1}\tau^3} +
\sup_{\tau \in R_2 \bigcup R_3}
    \exp\left(-\frac{1}{4} \infoev \norm{\tau}_2^2 \right) +
\ordlog{N^{-1}}
%
\\={}&
O(N^{-2}) +
\exp\left(-\frac{1}{4} \infoev \delta_1^2 (\log N)^2 \right) +
\ordlog{N^{-1}}
\\={}& \ordlog{N^{-1}},
\end{align*}
%
where, in the penultimate line, we use \lemref{region_r1_gaussian_moments},
including the given lower bound on $\delta_1$.
%
\end{proof}
%
\end{lem}

%%%%%%%%%%%%%%%%%%%%%%%%%%%%%%%%%%%%%%%%%%%%%%%%%%%%%%%%%%%%



%%%%%%%%%%%%%%%%%%%%%%%%%%%%%%%%%%%%%%%%%%
%%%%%%%%%%%%%%%%%%%%%%%%%%%%%%%%%%%%%%%%%%
%%%%%%%%%%%%%%%%%%%%%%%%%%%%%%%%%%%%%%%%%%




In \lemref{region_r1_final_bound} to follow, we integrate the result of
\lemref{region_r1_loglik_expansion} using \lemref{region_r1_gaussian_moments} to
produce our final results for region $R_1$. Recall from
\defref{loglik_details_def} that $\normhat = \normalizer^{-1}
\exp\left(\frac{1}{2} \infohat \tau^2 \right)$ is the density of the Gaussian
distribution on $\tau$ with mean zero and covariance $\infohat^{-1} =
(-\likhatk{2}(\thetahat))^{-1}$.

Note that the statement of \lemref{region_r1_final_bound} combines vector-valued
quantities like $\phigrad{2}(\thetahat, \xvec_s) \expect{\normhat}{\tau^2}$ with
scalar--valued quantities like $\norm{\phigrad{2}(\thetahat, \xvec_s)}_2$.  This is
for notational convenience since the scalar-valued terms will simply become part
of the residual, and their component values do not matter.  Formally, the
equation can be taken to hold componentwise.
%
\begin{lem}\lemlabel{region_r1_final_bound}
%
Let \assuref{core_bclt_assu} hold.  When $\delta_1$ is as given in
\lemref{region_r2}, then, for any $\xvec$ such that $\regevent$
occurs,
%
\begin{align*}
%
\MoveEqLeft
I_1(\phi) =
    \phi(\thetahat, \xvec_s)
    \int_{R_1} \exp\left(
        N \likhat(\theta) - N \likhat(\thetahat)
    \right) d\tau +
\nonumber\\ &\quad
    N^{-1} \normalizer \left(
        \frac{1}{2} \phigrad{2}(\thetahat, \xvec_s) \expect{\normhat}{\tau^2} +
        \frac{1}{6} \likhatk{3}(\thetahat) \phigrad{1}(\thetahat, \xvec_s)
            \expect{\normhat}{\tau^4} \right) +
\nonumber\\ &\quad
    \ordlog{N^{-2}} \bigg(
        1 +
        \norm{\phigrad{1}(\thetahat, \xvec_s)}_2 +
        \norm{\phigrad{2}(\thetahat, \xvec_s)}_2 +
        \expect{\normresid{\infoev / 2}{R_1}(\tau)}{
            \norm{\tresid{3}{\phi(\cdot, \xvec_s), \theta, \thetahat}}_2}
    \bigg).
%
\end{align*}
%
\begin{proof}
%
The result will follow by Taylor expanding $\phi$, combining with
\lemref{region_r1_loglik_expansion}, and then integrating using the fact that
Gaussian integrals in region $R_1$ are approximately equal to Gaussian integrals
over the whole domain, as proved in \lemref{region_r1_gaussian_moments}.


%%%%%%%%%%%%%%%%%

We first Taylor expand $\phi(\theta, \xvec_s)$ around $\thetahat$ and substitute
$\tau = \sqrt{N}(\theta - \thetahat)$.
%
\begin{align}
\phi(\theta, \xvec_s) ={}&
    \phi(\thetahat, \xvec_s) +
    N^{-1/2} \phigrad{1}(\thetahat, \xvec_s) \tau + \nonumber \\
&
    N^{-1} \frac{1}{2} \phigrad{2}(\thetahat, \xvec_s) \tau^2 +
    N^{-3/2} \tresid{3}{\phi(\cdot, \xvec_s), \theta, \thetahat} \tau^3.
    \eqlabel{r1_phi_expansion}
\end{align}


%%%%%%%%%%%%%%%%%%%%%%%%

We will now combine \eqref{r1_phi_expansion} with
\lemref{region_r1_loglik_expansion}.
%
\begin{align}
\MoveEqLeft
\exp\left(N \likhat(\theta) - N \likhat(\thetahat)\right)
\left(
    \phi(\theta, \xvec_s) - \phi(\thetahat, \xvec_s)
\right) = \nonumber \\
%%%%%%%%%%%%%%%%%%%%%%%
\MoveEqLeft
\exp\left(
    \frac{1}{2} \infohat \tau^2
\right) \times \nonumber\\
&
\left(
    N^{-1/2} \phigrad{1}(\thetahat, \xvec_s) \tau +
    N^{-1} \frac{1}{2} \phigrad{2}(\thetahat, \xvec_s) \tau^2 +
    N^{-3/2} \tresid{3}{\phi(\cdot, \xvec_s), \theta, \thetahat} \tau^3
\right)  \times \nonumber\\
&
\left(
    1 +
    N^{-1/2} \frac{1}{6} \likhatk{3}(\thetahat) \tau^3 +
    \ordlog{N^{-1}}
\right) = \nonumber\\
%%%%%%%%%%%%%%%%%%%%%%%
\MoveEqLeft
\exp\left(
    \frac{1}{2} \infohat \tau^2
\right) \times \nonumber\\
&
\Bigg(
    N^{-1/2} \phigrad{1}(\thetahat, \xvec_s) \tau +
\nonumber\\ &\quad
    N^{-1} \left(
        \frac{1}{2} \phigrad{2}(\thetahat, \xvec_s) \tau^2 +
        \frac{1}{6} \likhatk{3}(\thetahat) \phigrad{1}(\thetahat, \xvec_s) \tau^4 \right) +
\nonumber\\ &\quad
    N^{-3/2} \left(
        \tresid{3}{\phi(\cdot, \xvec_s), \theta, \thetahat} \tau^3 +
        \frac{1}{12} \likhatk{3}(\thetahat)  \phigrad{2}(\thetahat, \xvec_s) \tau^5
    \right) +
\nonumber\\ &\quad
    \ordlog{N^{-3/2}} \bigg(
        \phigrad{1}(\thetahat, \xvec_s) \tau +
% \nonumber\\ &\quad\quad\quad
        N^{-1/2}  \frac{1}{2} \phigrad{2}(\thetahat, \xvec_s) \tau^2 +
        N^{-1} \tresid{3}{\phi(\cdot, \xvec_s), \theta, \thetahat} \tau^3
    \bigg)
\Bigg) = \nonumber\\
%%%%%%%%%%%%%%%%%%%%%%%
\MoveEqLeft
\exp\left(
    \frac{1}{2} \infohat \tau^2
\right) \times \nonumber\\
&
\Bigg(
    N^{-1/2} \phigrad{1}(\thetahat, \xvec_s) \tau +
% \nonumber\\ &\quad
    N^{-1} \left(
        \frac{1}{2} \phigrad{2}(\thetahat, \xvec_s) \tau^2 +
        \frac{1}{6} \likhatk{3}(\thetahat) \phigrad{1}(\thetahat, \xvec_s) \tau^4 \right) +
\nonumber\\ &\quad
    \ordlog{N^{-3/2}} \bigg(
        \phigrad{1}(\thetahat, \xvec_s) \tau +
        \left( \ord{1} \tau^5 + \ord{N^{-1/2}}  \tau^2\right)
            \phigrad{2}(\thetahat, \xvec_s) +
\nonumber\\ &\quad\quad\quad\quad\quad\quad\quad
        \ord{1} \tresid{3}{\phi(\cdot, \xvec_s), \theta, \thetahat} \tau^3
    \bigg)
\Bigg).
\eqlabel{r1_product_expansion1}
\end{align}

In the final line of \eqref{r1_product_expansion1} we absorbed some higher-order
quantities into the lower-order $\ordlog{N^{-3/2}}$ term for convenience of
expression, and used the fact that, when $\regevent$ occurs, for $k=1,2,3$,
$\norm{\likhatk{k}(\thetahat)}_2 = \ord{1}$ when $\regevent$ occurs, since
%
\begin{align*}
%
\norm{\likhatk{k}(\thetahat) - \likk{k}(\thetahat)}_2 \le&
    \sup_{\theta \in R_1} \norm{\likhatk{k}(\theta) - \likk{k}(\theta)}_2 \le
    \epsilon_U,
%
\end{align*}
%
so
%
\begin{align*}
%
\norm{\likhatk{k}(\thetahat)}_2 \le&
    \norm{\likhatk{k}(\thetahat) - \likk{k}(\thetahat)}_2 +
    \norm{\likk{k}(\thetahat) - \likk{k}(\thetatrue)}_2 +
    \norm{\likk{k}(\thetatrue)}_2 \\
\le&
    \epsilon_U +
    \sup_{\theta \in R_1} \norm{\likk{k}(\theta) - \likk{k}(\thetatrue)}_2 + \norm{\likk{k}(\thetatrue)}_2
= \ord{1},
%
\end{align*}
%
by continuity of $\likk{k}(\theta)$.

Next, we integrate both sides of \eqref{r1_product_expansion1} using
\lemref{region_r1_gaussian_moments} to replace Gaussian integrals over $R_1$
with Gaussian moments over the whole domain, up to an order $N^{-3}$ which can
be absorbed in the $\ordlog{N^{-2}}$ error term.
%
\begin{align*}
%
I_1(\phi) ={}&
\int_{R_1} \phi(\theta, \xvec_s)
   \exp\left(N \likhat(\theta) - N \likhat(\thetahat)\right) d \tau
\\
={}&
\phi(\thetahat, \xvec_s)
\int_{R_1} \exp\left(
    N \likhat(\theta) - N \likhat(\thetahat)
\right) d\tau +
\\{}&
N^{-1} \left(
    \frac{1}{2} \phigrad{2}(\thetahat, \xvec_s)
        \expect{\normhat}{\tau^2} +
    \frac{1}{6} \likhatk{3}(\thetahat) \phigrad{1}(\thetahat, \xvec_s)
        \expect{\normhat}{\tau^4} \right) +
\nonumber\\ &
    \ordlog{N^{-2}} \bigg(
        1 +
        \norm{\phigrad{1}(\thetahat, \xvec_s)}_2 +
        \norm{\phigrad{2}(\thetahat, \xvec_s)}_2 +
\nonumber\\ &
    \hspace{5em}
        \int_{R_1}
            \exp\left(-\frac{1}{2} \infohat \tau^2 \right)
            \tresid{3}{\phi(\cdot, \xvec_s), \theta, \thetahat}
            d\tau
\bigg).
%
\end{align*}

Finally, since $\tresid{3}{\phi(\cdot, \xvec_s), \theta, \thetahat}$ depends
explicitly on $\tau$, we must control this term in a more convenient form.  We
will argue as in \lemref{region_r1_gaussian_moments}.   Recall that, when
$\regevent$ occurs, $\infoevhat \ge \frac{\infoev}{2}$, so that $\normalizer =
O(1)$ and
%
\begin{align*}
%
\MoveEqLeft
\norm{\normalizer \int_{R_1} \exp\left(-\frac{1}{2} \infohat \tau^2 \right)
    \tresid{3}{\phi(\cdot, \xvec_s), \theta, \thetahat} d\tau }_2
\\\le&
\normalizer \int_{R_1} \exp\left(-\frac{1}{2} \infohat \tau^2 \right)
    \norm{\tresid{3}{\phi(\cdot, \xvec_s), \theta, \thetahat}}_2  d\tau
\\ \le&
\normalizer \int_{R_1} \exp\left(-\frac{\infoev}{4} \norm{\tau}_2^2 \right)
    \norm{\tresid{3}{\phi(\cdot, \xvec_s), \theta, \thetahat}}_2  d\tau.
% \\ \le&
% \normalizer \int_{R_1 \bigcup R_2}
%     \exp\left(-\frac{\infoev}{4} \norm{\tau}_2^2 \right)
%         \norm{\tresid{3}{\phi(\cdot, \xvec_s), \theta, \thetahat}}_2  d\tau.
%
\end{align*}
%
The result follows by multiplying by the $O(1)$ normalizer of
the $\normresid{\infoev / 2}{R_1}$ distribution.
%
\end{proof}
%
\end{lem}
